\chapter{Decimation, Interpolation, and Resampling}

In Chapter~21, we estimated power spectra across frequency grids, implicitly
assuming a single, fixed sampling clock throughout the analysis. But a sample rate is
a choice, not an immutable law. When an analog-to-digital converter hands you a
stream of numbers and says ``these arrived at $48{,}000$ per second,'' it has
fixed a clock that every later interpretation depends on. But the world rarely
cooperates by giving you exactly the rate you want. Your accelerometer streams
at $104$ Hz, your gyroscope at $208$ Hz, your microphone at $44{,}100$ Hz, and
the gesture model you trained expects everything resampled to a tidy $100$ Hz
grid. Somewhere, somebody has to change a sample rate.

Changing a sample rate sounds like bookkeeping --- keep every other sample, or
squeeze new samples in between --- and if you do it that naively you will
silently corrupt your data. This chapter is about doing it correctly, and more
importantly about \emph{seeing} why the naive versions fail. Everything here is
a direct consequence of two ideas you already own: the spectrum of a sampled
signal is periodic (Chapter~18), and filtering is shaping that
spectrum (Chapters~14 and~16). Decimation, interpolation, and rational
resampling are just those two ideas applied with care.

\section{What it means to change a sample rate}
\label{ch22:meaning}

Let us be precise about the object we are manipulating. A discrete signal
$x[n]$ is a sequence of numbers. By itself it carries no time information at
all --- $x[0], x[1], x[2], \dots$ is just a list. The sample rate $f_s$ (in
samples per second, Hz) is the extra fact that tells us each index $n$
corresponds to physical time $t = n/f_s = n\,T$, where $T = 1/f_s$ is the
sampling period.

To ``change the sample rate'' means: produce a \emph{new} sequence $y[m]$,
sampled at a new rate $f_s'$, that represents the \emph{same underlying
continuous signal} as $x[n]$ did. The phrase ``same underlying continuous
signal'' is our guiding constraint. We are not free to invent numbers; we are obligated
to land on the values that the original analog waveform would have produced had
we sampled it at the new rate in the first place.

\begin{intuitionbox}
Think of $x[n]$ as marks you painted on a moving conveyor belt at regular
spacing. Decimation is stepping back so the marks look closer together (fewer
marks per meter of belt); interpolation is painting extra marks in between. In
both cases the belt --- the real signal --- never moved. You are only changing
how densely you describe it. The trap is that the marks already \emph{hide}
information about what happened between them, and a careless redescription can
make the belt appear to carry a pattern that was never there.
\end{intuitionbox}

Three operations cover essentially all cases:
\begin{itemize}
  \item \textbf{Decimation} (downsampling): lower the rate by an integer factor
  $M$, so $f_s' = f_s/M$. Fewer samples.
  \item \textbf{Interpolation} (upsampling): raise the rate by an integer
  factor $L$, so $f_s' = L\,f_s$. More samples.
  \item \textbf{Rational resampling}: change the rate by a ratio $L/M$, e.g.
  $44100 \to 48000$ is $\times 160/147$. Combine the two.
\end{itemize}
A truly arbitrary (irrational) ratio is handled by approximating it with a
rational one, or by computing samples on demand with a continuous-time
interpolation kernel; we will meet both.

\begin{focusbox}
The single most important rule of this chapter, stated up front so you can hold
it while we build the reasons:
\begin{center}
\emph{Before you throw samples away, remove the frequencies that the lower rate
cannot represent. After you add samples, remove the spectral images that the
upsampling created.}
\end{center}
Anti-alias on the way down, anti-image on the way up. Everything else is
detail.
\end{focusbox}

\section{Decimation: the geometry of throwing samples away}
\label{ch22:decim}

Start with the most innocent-looking operation imaginable. To halve the sample
rate, keep every second sample and discard the rest:
\begin{equation}
  y[m] = x[mM], \qquad m \in \Z.
  \label{ch22:downdef}
\end{equation}
For $M=2$: $y = (x[0], x[2], x[4], x[6], \dots)$. This is called
\emph{downsampling} or \emph{compressing}, and the block-diagram symbol is a
downward arrow with the factor: $\downarrow M$.

In the time domain nothing could be simpler --- we are reading the same list with a
stride of $M$. In the frequency domain, however, throwing away $M-1$ out of every
$M$ samples stretches the spectrum by $M$ and folds high frequencies straight onto low ones.

\subsection{What downsampling does to the spectrum}

Recall the central fact about sampling: a signal sampled at rate $f_s$ has a
spectrum that repeats every $f_s$. The discrete-time Fourier transform (DTFT)
is periodic with period $2\pi$ in normalized frequency $\omega = 2\pi f/f_s$.
When we downsample by $M$, the new rate is $f_s/M$, so the new spectrum must be
periodic with the \emph{smaller} period $f_s/M$. The only way to take a
spectrum periodic on $f_s$ and make it periodic on $f_s/M$ is to overlay $M$
shifted copies on top of each other. Formally,
\begin{equation}
  Y(e^{\jj\omega}) = \frac{1}{M}\sum_{k=0}^{M-1}
    X\!\left(e^{\jj(\omega - 2\pi k)/M}\right).
  \label{ch22:downspec}
\end{equation}
Read this slowly. The downsampled spectrum is the original spectrum
\emph{stretched} by $M$ (the $\omega/M$), \emph{scaled} in height by $1/M$, and
then \emph{summed over $M$ shifted versions} (the $2\pi k$ shifts). The stretch
is benign. The sum is where the danger lives: those $M$ copies can land on top
of one another. When they overlap, frequencies that were distinct in $x$ add
together in $y$ and become inseparable. That overlap is \emph{aliasing}, the
exact same phenomenon as undersampling an analog signal --- here it is induced
in the digital domain by the act of decimation.

\begin{whybox}
Why must the copies overlap unless we prevent it? After stretching by $M$, the
content that originally lived in $[0, \pi]$ now occupies $[0, M\pi]$ --- it is
$M$ times too wide to fit in one period of length $2\pi$. The replication then
folds the excess back in. The only way to avoid the fold is to make sure the
original content was narrow enough \emph{before} stretching that, even after
stretching by $M$, it still fits inside a single $2\pi$ window. That means the
original signal must occupy only $\abs{\omega} < \pi/M$, i.e. frequencies below
$f_s/(2M)$ --- precisely the new Nyquist limit.
\end{whybox}

\begin{figure}[H]
\centering
\begin{tikzpicture}
\begin{axis}[
  name=top,
  width=12cm, height=3.8cm,
  title={Original spectrum before decimation: energy spills past new Nyquist},
  title style={font=\small\bfseries},
  xlabel={frequency $\omega$ at original rate (rad/sample)},
  ylabel={$\abs{X(e^{\jj\omega})}$},
  xmin=-3.3, xmax=3.3, ymin=0, ymax=1.35,
  xtick={-3.14159,-1.5708,0,1.5708,3.14159},
  xticklabels={$-\pi$,$-\frac{\pi}{2}$,$0$,$\frac{\pi}{2}$,$\pi$},
  ytick={0,0.5,1.0}, axis on top,
]
% Baseband signal
\addplot[sigblue,fill=fillsig,fill opacity=0.45,domain=-1.5708:1.5708,samples=50,thick]
  {1-abs(x)/2.3562} \closedcycle;
% Excess spilling past pi/2
\addplot[trapred,fill=trapred,fill opacity=0.35,domain=1.5708:2.3562,samples=30,thick]
  {1-abs(x)/2.3562} \closedcycle;
\addplot[trapred,fill=trapred,fill opacity=0.35,domain=-2.3562:-1.5708,samples=30,thick]
  {1-abs(x)/2.3562} \closedcycle;
\draw[trapred,very thick,dashed] (axis cs:-1.5708,0) -- (axis cs:-1.5708,1.25);
\draw[trapred,very thick,dashed] (axis cs: 1.5708,0) -- (axis cs: 1.5708,1.25);
\node[trapred,font=\footnotesize] at (axis cs:1.95,0.7) {excess};
\node[trapred,font=\footnotesize] at (axis cs:-1.95,0.7) {excess};
\node[font=\footnotesize,anchor=south] at (axis cs:0,1.05) {passband to keep};
\end{axis}

\begin{axis}[
  at={($(top.below south west)+(0,-0.85cm)$)},
  anchor=north west,
  width=12cm, height=3.8cm,
  title={Downsampled spectrum: stretch by $M=2$ causes spectral folding and aliasing},
  title style={font=\small\bfseries},
  xlabel={frequency $\omega'$ at downsampled rate (rad/sample)},
  ylabel={$\abs{Y(e^{\jj\omega'})}$},
  xmin=-3.3, xmax=3.3, ymin=0, ymax=0.75,
  xtick={-3.14159,-1.5708,0,1.5708,3.14159},
  xticklabels={$-\pi$,$-\frac{\pi}{2}$,$0$,$\frac{\pi}{2}$,$\pi$},
  ytick={0,0.25,0.5}, axis on top,
]
% Stretched baseband copy (scaled by 1/2)
\addplot[sigblue,dashed,thick,domain=-3.14159:3.14159,samples=70]
  {0.5*max(0, 1 - abs(x)/(2*2.3562))};
% Shifted replica from +2pi
\addplot[trapred,dashed,thick,domain=-3.14159:3.14159,samples=70]
  {0.5*max(0, 1 - abs(x - 6.28318)/(2*2.3562)) + 0.5*max(0, 1 - abs(x + 6.28318)/(2*2.3562))};
% Folded sum (aliased spectrum)
\addplot[structurepurple,very thick,domain=-3.14159:3.14159,samples=120]
  {0.5*max(0, 1 - abs(x)/(4.7124)) + 0.5*max(0, 1 - abs(x - 6.28318)/(4.7124)) + 0.5*max(0, 1 - abs(x + 6.28318)/(4.7124))};
\node[structurepurple,font=\footnotesize,anchor=west] at (axis cs:1.2,0.55) {Folded sum $Y(e^{\jj\omega'})$};
\node[sigblue,font=\scriptsize,anchor=south] at (axis cs:0,0.51) {baseband};
\node[trapred,font=\scriptsize,anchor=south] at (axis cs:2.0,0.10) {alias replica};
\end{axis}
\end{tikzpicture}
\caption{Downsampling spectral folding. Top: The original spectrum $X(e^{\jj\omega})$ extends beyond the new Nyquist limit $\pm \pi/M = \pm \pi/2$ (red excess). Bottom: Downsampling stretches the spectrum by $M=2$ and superimposes shifted copies spaced by $2\pi$. The shifted replica folds directly into the baseband $[-\pi, \pi]$, creating unrecoverable aliasing distortion (purple sum) unless removed by a lowpass filter prior to downsampling.}
\label{ch22:fig-predec}
\end{figure}

\subsection{The fix: filter, then decimate}

The remedy is forced on us by the geometry. Before stride-reading, run $x[n]$
through a lowpass filter that removes everything above the new Nyquist frequency
$f_s/(2M)$, i.e. above $\omega_c = \pi/M$. With the offending frequencies gone,
the shifted copies in \eqref{ch22:downspec} no longer overlap, and the
decimation is information-preserving for the band we kept. The combined
operation --- lowpass then $\downarrow M$ --- is what most engineers actually
mean by ``decimation.''
\begin{equation}
  x[n] \;\longrightarrow\; \boxed{\text{LPF, cutoff } \pi/M}
  \;\longrightarrow\; \boxed{\downarrow M} \;\longrightarrow\; y[m].
  \label{ch22:decimchain}
\end{equation}

\begin{trapbox}
The classic beginner bug: ``I'll just take every other sample, it's faster.'' If
the signal has \emph{any} energy above the new Nyquist limit --- and almost all
real signals do, because of wideband noise alone --- that energy folds down and
contaminates your low band permanently. There is no later filter that can remove
it, because after aliasing the noise and the signal share the same frequency
bins. The damage is done at the moment of decimation. Filter first.
\end{trapbox}

\begin{example}[Decimating an IMU stream from $208$ to $52$ Hz]
\label{ch22:ex-imu}
An ICM-20948 gyroscope is configured for $208$ Hz output, but our gesture model
wants $52$ Hz, so $M = 4$. The new Nyquist is $26$ Hz. Vibration from a tapping
finger or a nearby motor can easily put energy at, say, $40$ Hz. If we naively
keep every 4th sample, that $40$ Hz tremor aliases: it lands at
$\abs{40 - 52} = 12$ Hz --- right in the band where deliberate wrist motions
live (a few Hz to $\sim 15$ Hz). It would masquerade as a real gesture. A
lowpass at $26$ Hz before decimating kills the tremor before it can fold.
\end{example}

\begin{happeningbox}
What is really happening when we ``decimate'' is that we admit our new, coarser
clock is incapable of describing fast wiggles, so we deliberately blur the
signal down to wiggles it \emph{can} describe, and only then re-sample it on the
coarse grid. Decimation is deliberate forgetting: we choose what to lose, rather
than letting aliasing scramble it unpredictably.
\end{happeningbox}

\subsection{You only need to compute the outputs you keep}

Here is the first efficiency insight, and it foreshadows polyphase. In the chain
\eqref{ch22:decimchain} the filter computes an output for \emph{every} input
sample, but the $\downarrow M$ immediately discards $M-1$ of every $M$ of them.
That is wasteful: we are doing the full convolution at the high rate only to
throw most of it away. Since the filter output at the discarded indices is never
used, we can simply \emph{not compute it}. If the lowpass has $N$ taps, the
direct chain costs $N$ multiplies per input sample; computing only the kept
outputs costs $N$ multiplies per \emph{output} sample, i.e. $N/M$ per input
sample --- an $M$-fold saving. We make this rigorous with polyphase
decomposition in Section~\ref{ch22:polyphase}.

\section{Interpolation: the geometry of inventing samples}
\label{ch22:interp}

Now go the other way. We want more samples than we have. The correct version
proceeds in two steps that mirror decimation in reverse.

\subsection{Step 1: zero-stuffing (the expander)}

To raise the rate by $L$, insert $L-1$ zeros after each sample:
\begin{equation}
  v[m] =
  \begin{cases}
    x[m/L], & m = 0, \pm L, \pm 2L, \dots\\
    0, & \text{otherwise.}
  \end{cases}
  \label{ch22:expand}
\end{equation}
The block symbol is $\uparrow L$. For $L=3$ and $x = (a, b, c)$ we get
$v = (a, 0, 0, b, 0, 0, c, 0, 0)$. We have not changed any existing value and we
have made the sequence $L$ times longer, so it now lives on a grid that is $L$
times finer. So far so harmless --- but spectrally something curious happens.

\subsection{What zero-stuffing does to the spectrum}

Zero-stuffing does \emph{not} change the spectrum's shape at all; it merely
relabels the frequency axis. Because we packed the same content into $L$ times as
many samples, the new (higher) rate compresses the frequency axis, and the
original spectrum, which lived once per $2\pi$, now appears $L$ times within one
new period:
\begin{equation}
  V(e^{\jj\omega}) = X(e^{\jj\omega L}).
  \label{ch22:expandspec}
\end{equation}
The DTFT of $v$ is just the DTFT of $x$ compressed by $L$ --- which means $L$
copies of the baseband spectrum now appear in $[-\pi, \pi]$. The extra $L-1$
copies are called \emph{spectral images}. They are not noise; they are exact
replicas of the signal sitting at higher frequencies, created purely by the act
of inserting zeros.

\begin{figure}[H]
\centering
\begin{tikzpicture}
\begin{axis}[
  width=12cm, height=4.8cm,
  title={Upsampling spectral images: zero-stuffing compresses frequency axis},
  title style={font=\small\bfseries},
  xlabel={normalized frequency $\omega$ at the new high rate (rad/sample)},
  ylabel={$\abs{V(e^{\jj\omega})}$},
  xmin=-3.3, xmax=3.3, ymin=0, ymax=1.5,
  xtick={-3.14159,-1.0472,0,1.0472,3.14159},
  xticklabels={$-\pi$,$-\frac{\pi}{3}$,$0$,$\frac{\pi}{3}$,$\pi$},
  ytick={0,0.5,1.0}, axis on top, every axis plot/.append style={thick},
]
\addplot[sigblue,fill=fillsig,fill opacity=0.55,domain=-0.5236:0.5236,samples=40]
  {1-abs(x)/0.5236} \closedcycle;
\addplot[specred,fill=specred,fill opacity=0.30,domain=1.5708:2.618,samples=40]
  {1-abs(x-2.0944)/0.5236} \closedcycle;
\addplot[specred,fill=specred,fill opacity=0.30,domain=-2.618:-1.5708,samples=40]
  {1-abs(x+2.0944)/0.5236} \closedcycle;
\draw[targetgreen,very thick,dashed] (axis cs:-1.0472,0) -- (axis cs:-1.0472,1.25)
  -- (axis cs:1.0472,1.25) -- (axis cs:1.0472,0);
\node[targetgreen,anchor=south,font=\small] at (axis cs:0,1.27)
  {anti-image LPF passes $\abs{\omega}<\pi/L$};
\node[specred,anchor=south,font=\footnotesize] at (axis cs:2.0944,1.03) {spectral image};
\node[specred,anchor=south,font=\footnotesize] at (axis cs:-2.0944,1.03) {spectral image};
\node[sigblue,anchor=south,font=\footnotesize] at (axis cs:0,1.03) {baseband};
\end{axis}
\end{tikzpicture}
\caption{After $\uparrow L$ with $L=3$. The baseband copy (blue) is the signal we
want; the red copies are spectral images created by zero-stuffing. The
interpolation lowpass (green dashed line) keeps only $\abs{\omega}<\pi/L$ and erases the
images, smoothly filling in the zeros.}
\label{ch22:fig-images}
\end{figure}

\subsection{Step 2: the anti-imaging (interpolation) filter}

The zeros we inserted are placeholders; we want the filter to replace them with
sensible interpolated values. In the frequency domain that is exactly the same as
removing the images: keep only the baseband copy $\abs{\omega} < \pi/L$ and zero
out the rest with a lowpass filter of cutoff $\pi/L$. So:
\begin{equation}
  x[n] \;\longrightarrow\; \boxed{\uparrow L} \;\longrightarrow\;
  \boxed{\text{LPF, cutoff } \pi/L,\ \text{gain } L} \;\longrightarrow\; y[m].
  \label{ch22:interpchain}
\end{equation}
The gain $L$ compensates for the energy spreading: a unit-height passband would
make the output $L$ times too quiet, because only $1/L$ of the samples carry all
the energy. With gain $L$ the interpolated samples have the right amplitude.

\begin{whybox}
Why is ``remove the images'' the same as ``fill in the zeros sensibly''? In the
time domain, convolving the zero-stuffed sequence with the filter's impulse
response $h$ places a copy of $h$, scaled by $x[k]$, at every original sample
position $kL$, and sums them. The zeros are filled by the overlapping skirts of
those kernels. If $h$ is an ideal lowpass (a sinc), the result is the unique
band-limited interpolation --- the same curve the Whittaker--Shannon formula
gives. ``Kill the images'' and ``draw the smooth band-limited curve through the
samples'' are literally the same operation viewed from two domains.
\end{whybox}

\subsection{Interpolation kernels you can picture}

The choice of $h$ is a choice of how you believe the signal behaves between
samples. Three kernels, in increasing quality and cost:
\begin{itemize}
  \item \textbf{Zero-order hold (nearest)}: $h$ is a rectangle of width $L$. Every
  inserted sample copies its left neighbor. Cheap, blocky; its frequency response
  is a wide $\sinc$ that passes a lot of image energy. Staircase output.
  \item \textbf{Linear interpolation}: $h$ is a triangle of width $2L$. Inserted
  samples sit on straight lines between originals. Better; frequency response is
  $\sinc^2$, so images are attenuated more. Still rolls off the top of the
  passband.
  \item \textbf{Bandlimited (windowed sinc)}: $h$ approximates the ideal
  brick-wall lowpass. Sharpest image rejection, flattest passband, highest cost.
  This is what high-quality audio resamplers use.
\end{itemize}

\begin{figure}[H]
\centering
\begin{tikzpicture}
\begin{axis}[
  width=12cm, height=4.4cm,
  xlabel={time (samples)}, ylabel={amplitude},
  xmin=-0.3, xmax=8.3, ymin=-0.1, ymax=1.25,
  xtick={0,2,4,6,8}, ytick=\empty,
  legend style={font=\small,at={(0.99,0.98)},anchor=north east},
  every axis plot/.append style={thick},
]
\addplot[only marks,mark=*,sigblue,mark size=2.4pt]
  coordinates {(0,0.2)(4,1.0)(8,0.3)};
\addlegendentry{originals ($\uparrow 4$)}
\addplot[quotientorange,const plot mark left] coordinates
  {(0,0.2)(4,1.0)(8,0.3)};
\addlegendentry{zero-order hold}
\addplot[targetgreen] coordinates {(0,0.2)(4,1.0)(8,0.3)};
\addlegendentry{linear}
\end{axis}
\end{tikzpicture}
\caption{Three originals (dots) upsampled by $L=4$. Zero-order hold (orange)
gives a staircase; linear (green) gives straight segments. A windowed-sinc filter
would give a smooth band-limited curve through the same dots.}
\label{ch22:fig-kernels}
\end{figure}

\begin{connectionbox}
The interpolation filter \emph{is} the discrete analogue of the
Whittaker--Shannon reconstruction integral from Chapter~18, and that
integral is itself an orthogonal projection: it expresses the band-limited signal in the
orthonormal basis of shifted sincs (Chapters~3 and~42).
Choosing a cheaper kernel is choosing a non-ideal, non-orthogonal basis --- you
trade reconstruction error for arithmetic. Every resampler is a statement about
which basis you can afford.
\end{connectionbox}

\section{Rational resampling: changing rate by $L/M$}
\label{ch22:rational}

Most real rate conversions are neither pure integers up nor down. To go from
$44.1$ kHz to $48$ kHz you multiply by $48000/44100 = 160/147$. The recipe is to
do both, \emph{in the right order}: upsample by $L$, then downsample by $M$, with
a single lowpass in the middle that serves \emph{both} as anti-image (for the
$\uparrow L$) and anti-alias (for the $\downarrow M$):
\begin{equation}
  x \;\to\; \boxed{\uparrow L} \;\to\;
  \boxed{\text{LPF, cutoff } \min\!\big(\tfrac{\pi}{L},\tfrac{\pi}{M}\big),\
  \text{gain } L} \;\to\; \boxed{\downarrow M} \;\to\; y.
  \label{ch22:ratchain}
\end{equation}

\begin{focusbox}
\textbf{Always upsample first.} If you downsample before upsampling, the
$\downarrow M$ may have already discarded frequencies that the final rate
\emph{could} have kept, and no later upsampling recovers them. Going up first,
then filtering, then down preserves the maximum bandwidth the output rate can
support. The single shared filter must protect both edges, so its cutoff is the
\emph{lower} of the two limits, $\min(\pi/L, \pi/M)$.
\end{focusbox}

\begin{example}[$44.1$ kHz to $48$ kHz]
\label{ch22:ex-audio}
Reduce $48000/44100$: dividing by $\gcd = 300$ gives $160/147$. So $L=160$,
$M=147$. We conceptually upsample by $160$ (to $7.056$ MHz!), lowpass at the
lower of $\pi/160$ and $\pi/147$, namely $\pi/160$, then downsample by $147$. We
never actually materialize the $7$ MHz stream --- polyphase (next section) lets
us compute only the $48000$ outputs we keep, directly from the $44100$ inputs.
\end{example}

\begin{example}[A messy ratio: $104 \to 100$ Hz]
\label{ch22:ex-imu2}
An accelerometer at $104$ Hz, model wants $100$ Hz. Ratio $100/104 = 25/26$, so
$L=25$, $M=26$. Tiny rate change, but the same machinery applies: upsample by
$25$, lowpass at $\min(\pi/25,\pi/26)=\pi/26$, downsample by $26$. Because $L$ and
$M$ are close, the filter does very little reshaping and a modest design suffices.
Contrast this with sloppy alternatives: dropping one sample every $26$
(sample-and-hold decimation) introduces a periodic glitch at $4$ Hz --- squarely
inside the gesture band --- whereas the proper filter spreads the tiny correction
smoothly.
\end{example}

\begin{trapbox}
Reduce the ratio before choosing $L$ and $M$. Using $48000/44100$ unreduced
($L=48000$, $M=44100$) asks for a $2.1$ GHz intermediate rate and a filter with
millions of taps for no benefit. The reduced form $160/147$ gives the identical
result with a tractable design. Always divide by $\gcd(L,M)$ first.
\end{trapbox}

\section{Polyphase: doing it without the waste}
\label{ch22:polyphase}

The chains above are correct but, taken literally, absurd: in
Example~\ref{ch22:ex-audio} they build a $7$ MHz signal that is $99\%$ zeros,
filter all of it, then throw away $146$ of every $147$ outputs. Polyphase
decomposition is the algebra that removes every wasted multiply. The result is
not an approximation --- it computes \emph{exactly} the same output, just without
touching zeros or computing samples we discard.

\subsection{The idea in one picture}

Consider interpolation by $L$. The zero-stuffed signal $v$ is mostly zeros, and
when we convolve it with the filter $h$, every multiply by one of those zeros
contributes nothing. At any given output index, only the taps of $h$ that happen
to land on the nonzero (original) samples do real work. As the output index
advances through one period of $L$, a \emph{different} subset of taps lights up
each time. So instead of one length-$N$ filter run at the high rate, we have $L$
smaller \emph{subfilters} (the polyphase components), each run at the low rate,
and we interleave their outputs.

\begin{intuitionbox}
Picture the filter's $N$ taps as a comb, and the nonzero samples as teeth spaced
$L$ apart. Slide the comb past the teeth. At each shift, only the comb positions
sitting exactly on a tooth matter; the rest hover over zeros. There are exactly
$L$ distinct alignments before the pattern repeats. Each alignment \emph{is} a
polyphase subfilter. Polyphase is simply the observation that you should evaluate
each alignment separately and skip the empty positions.
\end{intuitionbox}

\subsection{The decomposition, written out}

Split the filter's impulse response into $L$ interleaved phases. The $\rho$-th
polyphase component is
\begin{equation}
  h_\rho[n] = h[nL + \rho], \qquad \rho = 0, 1, \dots, L-1.
  \label{ch22:polydef}
\end{equation}
That is, $h_0$ takes taps $h[0], h[L], h[2L], \dots$; $h_1$ takes
$h[1], h[L+1], \dots$; and so on. Together the $L$ subfilters contain every tap
of $h$ exactly once. For interpolation, output phase $\rho$ is produced by
filtering the \emph{original} (not zero-stuffed) signal with $h_\rho$:
\begin{equation}
  y[mL + \rho] = \sum_k x[k]\,h_\rho[m-k] = (x \conv h_\rho)[m].
  \label{ch22:polyinterp}
\end{equation}
Each of the $L$ subfilters has about $N/L$ taps and runs at the input rate, so
the total work is $\sim N$ multiplies per input sample --- the same as one full
filter run, but with the $\times L$ rate inflation completely removed. The zeros
are never multiplied because they were never created.

For decimation by $M$ the dual statement holds: split the \emph{input} into $M$
phases, filter each by a subfilter, and \emph{sum} the results to form each kept
output. Again you compute only outputs you keep, at $\sim N/M$ multiplies per
input sample.

\begin{figure}[H]
\centering
\begin{tikzpicture}[
  font=\small,
  blk/.style={draw,rounded corners,minimum height=7mm,minimum width=15mm,fill=blue!4},
  >={Stealth[length=2.2mm]},
]
\node (x) at (0,0) {$x[n]$};
\node[blk] (h0) at (3,1.4) {$h_0$};
\node[blk] (h1) at (3,0)   {$h_1$};
\node[blk] (h2) at (3,-1.6){$h_{L-1}$};
\node at (3,-0.8) {$\vdots$};
\draw[->] (x) -- (h0); \draw[->] (x) -- (h1); \draw[->] (x) -- (h2);
\node[draw,rounded corners,minimum height=26mm,minimum width=15mm,fill=green!5]
  (cm) at (6.6,0) {\shortstack{commutator\\(interleave)}};
\draw[->] (h0) -- (cm.west|-h0);
\draw[->] (h1) -- (cm.west|-h1);
\draw[->] (h2) -- (cm.west|-h2);
\node (y) at (9.7,0) {$y[m]$};
\draw[->] (cm) -- (y);
\end{tikzpicture}
\caption{Polyphase interpolator. The input drives $L$ short subfilters
$h_0,\dots,h_{L-1}$ in parallel at the low rate; a commutator interleaves their
outputs to form the high-rate stream. No zeros are ever multiplied.}
\label{ch22:fig-polyphase}
\end{figure}

\begin{whybox}
Why is this exact and not an approximation? Because zero-stuffing followed by
convolution is a \emph{linear} operation, and \eqref{ch22:polyinterp} is just
that same linear map re-grouped. Multiplying by zero and adding the result is
algebraically removable without changing the sum. Polyphase reorders the
arithmetic so that the removable operations are never performed. The Noble
identities (which let you move a filter through a $\downarrow M$ or $\uparrow L$
block) are the formal license for this reordering; the payoff is an order-$M$ or
order-$L$ speedup with bit-for-bit identical output.
\end{whybox}

\begin{example}[Tap count saved]
\label{ch22:ex-taps}
A $48$-tap interpolation filter at $L=4$ run the naive way costs $48$ multiplies
for \emph{each} high-rate output, i.e. $192$ per input sample. Polyphase splits it
into four $12$-tap subfilters; producing all four output phases costs
$4 \times 12 = 48$ multiplies per input sample --- a $4\times$ saving, exactly the
upsampling factor.
\end{example}

\subsection{A minimal polyphase interpolator in code}

\begin{lstlisting}[language=Python,caption={Polyphase upsample by L; no zeros are multiplied.},captionpos=b]
import numpy as np

def polyphase_upsample(x, h, L):
    """Upsample x by integer L using prototype lowpass h (length ~ N*L)."""
    h = np.asarray(h, float)
    pad = (-len(h)) % L                  # pad so length is a multiple of L
    hh = np.concatenate([h, np.zeros(pad)])
    phases = hh.reshape(-1, L).T         # phases[rho] = h[rho], h[rho+L], ...
    outs = [np.convolve(x, phases[rho]) for rho in range(L)]
    y = np.empty(L * len(outs[0]))
    for rho in range(L):
        y[rho::L] = outs[rho]            # interleave: y[mL+rho] = outs[rho][m]
    return L * y                         # gain L compensates energy spreading
\end{lstlisting}

For everyday work you would not hand-roll this; SciPy's \texttt{resample\_poly}
does rational $L/M$ resampling with an internally designed polyphase filter, and
\texttt{scipy.signal.decimate} handles integer-$M$ decimation with the
anti-alias filter built in. Writing the loop out explicitly makes clear that the
library is doing nothing exotic --- just the interleave-filter-and-skip pipeline we drew.

\section{Arbitrary and time-varying rates}
\label{ch22:arbitrary}

Sometimes the ratio is not nicely rational, or it drifts: two sensors with
independent crystal oscillators are never \emph{exactly} at their nominal rates,
so a stream nominally at $100$ Hz might really be $100.012$ Hz and wandering with
temperature. Two practical tools:
\begin{itemize}
  \item \textbf{Farrow / fractional-delay interpolation.} Fit a low-order
  polynomial (often cubic) through a few neighboring samples and evaluate it at
  the desired fractional position $\mu \in [0,1)$ between input samples. The same
  coefficient structure works for \emph{any} $\mu$, so the output instant can be
  placed anywhere --- ideal when the resampling phase advances by a
  non-integer step each output. This is the standard engine of software-defined radio
  symbol timing recovery.
  \item \textbf{Rational approximation.} Pick $L/M$ close enough to the true
  ratio that the residual error is below your tolerance, then use polyphase. For
  $100.012/100$ a ratio like $25003/25000$ is overkill-accurate; a coarser one
  is usually fine.
\end{itemize}

\begin{labbox}
Generate a $7$ Hz sinusoid sampled at $200$ Hz. (1) Decimate to $50$ Hz the wrong
way (\texttt{x[::4]}) and the right way (\texttt{scipy.signal.decimate(x,4)});
overlay their spectra. (2) Add a $90$ Hz tone before decimating and watch where it
lands in each case --- the naive version aliases it to $90-4\cdot 25 = -10$ Hz,
i.e. a phantom $10$ Hz line that the proper filter removes. Seeing that phantom
appear demonstrates the foundational multirate danger in a single plot.
\end{labbox}

\section{Why this matters for fusing sensors at different rates}
\label{ch22:fusion}

\begin{goalbox}
The wrist controller does not have one signal --- it has a small zoo of them,
each on its own clock. A typical build streams:
\begin{itemize}
  \item accelerometer + gyroscope (IMU) at, say, $208$ Hz;
  \item magnetometer at a slow $50$ Hz (mag ODRs are usually low);
  \item if present, an EMG or force-myography (FMG) front end at $1000$+ Hz,
  because muscle activity carries energy up to several hundred Hz.
\end{itemize}
The sensor-fusion filter (orientation estimator) and the gesture classifier both
want a \emph{single} synchronous frame rate --- often $100$ Hz, a comfortable
grid for wrist motion whose deliberate content sits below $\sim 15$ Hz. So every
stream must be resampled onto that common $100$ Hz grid before fusion. The IMU is
\emph{decimated} (anti-alias at $50$ Hz, then down); the magnetometer is
\emph{interpolated} up; the EMG is heavily decimated, but only \emph{after} you
have extracted its high-frequency features (envelope, RMS, band powers), because
those features are what you actually fuse at $100$ Hz --- you do not feed raw
$1$ kHz EMG into a $100$ Hz model and expect its fast content to survive.

Get the resampling wrong and the fusion silently degrades: an un-anti-aliased IMU
decimation folds motor/tremor vibration into the gesture band and the classifier
learns to fire on it; a magnetometer interpolated with a staircase hold injects
step glitches that the orientation filter reads as sudden rotations. The boring,
unglamorous resampling stage is where a lot of ``my gesture detector is twitchy''
bugs actually live.
\end{goalbox}

\begin{trapbox}
Resampling moves samples onto a common rate; it does \emph{not} align them in
absolute time. Two sensors started at different instants, or read over a bus with
different latencies, are time-\emph{offset} even at the same rate. You must also
estimate and correct that offset (a fixed delay, or a slowly varying one if the
clocks drift) --- typically by cross-correlating a shared motion event, or by
timestamping every sample at acquisition. Same-rate is necessary but not
sufficient for ``synchronized.''
\end{trapbox}

\begin{aestheticbox}
Be clear-eyed about what resampling can and cannot buy you at the wrist.
Resampling fixes \emph{rate} mismatches; it adds no information. A wrist-worn IMU sampled and
resampled to any rate still cannot recover continuous individual finger-joint
angles --- the wrist simply does not move enough, per-finger, for the IMU to see
it. The IMU gives you coarse wrist orientation and discrete gross gestures.
Finger-level discrimination needs EMG or FMG on the forearm (and even then it is
classification of a handful of states, not a continuous joint readout); true
continuous finger pose needs sensors on the fingers themselves, such as
instrumented rings. Clean resampling makes whatever signal you \emph{do} have
usable across sensors; it cannot conjure a degree of freedom the sensor never
measured.
\end{aestheticbox}

\begin{buildbox}
A pragmatic on-device pipeline for a Cortex-M class MCU (e.g. the nRF52840 or an
ESP32 reading the IMU over SPI):
\begin{enumerate}
  \item Read each sensor in its own ISR/queue at its native ODR; timestamp every
  sample from a single monotonic clock.
  \item Per stream, run a short anti-alias FIR (a $16$--$32$-tap polyphase
  decimator is plenty for IMU $208\to104\to52$ in stages --- cascading two
  $\downarrow 2$ stages is cheaper than one $\downarrow 4$, because each stage's
  filter only has to be sharp relative to its own band).
  \item Resample slow streams up to the common grid with linear interpolation if
  CPU is tight; the magnetometer changes slowly enough that linear is usually
  fine. Reserve windowed-sinc for anything whose passband you care about near
  Nyquist.
  \item Emit one fused frame per output tick; feed that to the orientation filter
  and the classifier.
\end{enumerate}
Cascaded half-band filters are the standard trick: a half-band FIR has nearly
half its taps equal to zero, so a $\downarrow 2$ half-band stage is remarkably
cheap, and several in series give large decimation factors at minimal cost. CIC
(cascaded integrator-comb) filters are an even cheaper, multiplier-free option
for very large integer factors --- common in the decimation stage of
sigma-delta ADCs --- at the price of a droopy passband you must later
compensate.
\end{buildbox}

\begin{applicationbox}
Rational resampling is everywhere once you look:
\begin{itemize}
  \item \textbf{Audio.} Every time a $44.1$ kHz CD track plays through a $48$ kHz
  sound card, a $160/147$ polyphase resampler runs. Sample-rate converters are a
  whole product category precisely because doing it transparently is hard.
  \item \textbf{Software radio.} A wideband receiver samples at tens of MHz, then
  decimates hard to bring one narrow channel down to a manageable baseband rate
  --- multi-stage polyphase/CIC decimation is the front end of essentially every
  SDR.
  \item \textbf{Biomedical.} ECG/EEG recorded at one rate are routinely resampled
  to a study's standard rate so records from different machines can be compared;
  done carelessly it shifts QRS timing.
  \item \textbf{Motion / wearables.} Exactly our case: aligning IMU, mag, and
  EMG/FMG onto one grid for fusion.
\end{itemize}
\end{applicationbox}

\begin{historybox}
The polyphase viewpoint and the multirate ``Noble identities'' were crystallized
in the 1970s--80s, with Maurice Bellanger's work on polyphase networks and the
later synthesis in P.~P.~Vaidyanathan's \emph{Multirate Systems and Filter Banks}
(1993), which became the standard reference. The same multirate algebra underlies
filter banks and, later, wavelets (Chapter~24) --- decimation and
interpolation are the atoms from which multiresolution analysis is built. What
looks like a humble rate-changing trick is the same machinery that powers JPEG-
2000 and MP3 subband coding.
\end{historybox}

\section{Putting the pieces together}
\label{ch22:summary}

Strip away the algebra and the chapter is two sentences. To go \emph{down} in
rate, first lowpass to the new Nyquist so nothing folds, then keep every $M$-th
sample. To go \emph{up}, insert $L-1$ zeros and then lowpass to erase the images
the zeros created. Rational changes do both with one shared filter, upsampling
first; polyphase makes all of it cheap by never multiplying a zero or computing a
sample you will discard. And in a real multi-sensor build, this is the quiet
stage that decides whether your fused signal is clean or full of phantoms.

Resampling lets us adapt sample rates between different stages of a pipeline.
Rate conversion shifts the spacing of spectral replicas, but what if we want to translate
a signal's spectrum to an entirely new center frequency, or transmit two independent signals
over the same frequency band? In Chapter~23 (\emph{Modulation and I/Q Signals}), we explore
amplitude modulation, frequency translation, and the complex quadrature representation
used across digital communications, radar, and sensor front-ends.

\section{Exercises}
\label{ch22:exercises}

\begin{exercise}
A signal occupies $\abs{\omega} < 0.6\pi$. You decimate by $M=2$ without
filtering. Sketch the resulting spectrum from \eqref{ch22:downspec} and identify
the band of frequencies that becomes corrupted by aliasing. What cutoff would the
anti-alias filter have needed?
\end{exercise}

\begin{exercise}
A tone at $f_0 = 4300$ Hz is sampled at $f_s = 16$ kHz, then decimated by $M=4$
\emph{without} an anti-alias filter. To what frequency does it alias in the
$4$ kHz output? Now repeat with a proper anti-alias filter and state what happens
to the tone instead.
\end{exercise}

\begin{exercise}
You upsample $x = (1, -1)$ by $L=3$ using zero-order hold, then by $L=3$ using
linear interpolation. Write out the two output sequences (ignore filter delay).
Which one has more high-frequency image energy, and why does the time-domain shape
tell you so?
\end{exercise}

\begin{exercise}
Show from \eqref{ch22:expandspec} that zero-stuffing alone (no filter) leaves the
\emph{values} at the original sample positions unchanged. Why, then, is the
interpolation filter still necessary --- what is wrong with the in-between zeros?
\end{exercise}

\begin{exercise}
Reduce the resampling ratios for the following and give $(L, M)$: (a)
$96000 \to 44100$; (b) $1000 \to 100$; (c) $208 \to 100$. For (c), state the
cutoff $\min(\pi/L, \pi/M)$ of the shared filter.
\end{exercise}

\begin{exercise}
A $64$-tap interpolation filter is used at $L=8$. Count the multiplies per input
sample for (a) the naive zero-stuff-then-filter implementation and (b) the
polyphase implementation. What is the ratio, and how does it relate to $L$?
\end{exercise}

\begin{exercise}
Explain, using \eqref{ch22:downspec} and \eqref{ch22:expandspec}, why rational
resampling must upsample \emph{before} downsampling. Construct a concrete $L/M$
and a signal where doing it in the wrong order loses recoverable bandwidth.
\end{exercise}

\begin{exercise}
A half-band lowpass filter has $h[n]=0$ for all even $n$ except the center tap.
Argue why this makes a $\downarrow 2$ polyphase decimator especially cheap, and
estimate the multiply savings versus a general filter of the same length.
\end{exercise}

\begin{exercise}
(Wrist tracker.) Your IMU runs at $208$ Hz and your magnetometer at $50$ Hz; both
must reach a common $100$ Hz grid. Give the resampling operation for each stream
(factors and filter type), and explain which stream is at risk of aliasing and
which is at risk of imaging.
\end{exercise}

\begin{exercise}
(What resampling cannot buy.) A teammate proposes resampling the wrist IMU to
$1000$ Hz so the model
can ``read individual finger movements.'' Explain in two or three sentences why
the higher rate cannot recover finger-joint angles, and what sensor change would
actually be required.
\end{exercise}

\begin{exercise}
(Code.) Modify the polyphase listing to perform decimation by $M$ instead of
interpolation: split the \emph{input} into $M$ phases, filter each by the
corresponding $h_\rho$, and sum. Verify on a test signal that your output matches
\texttt{scipy.signal.decimate} up to filter transients.
\end{exercise}

\begin{exercise}
Two nominally-$100$ Hz sensors actually run at $100.00$ and $100.05$ Hz. Over a
$60$-second capture, how many samples do they drift apart? Describe a resampling
strategy (rational approximation or fractional-delay) that keeps them aligned,
and state what additional information you need beyond the sample values.
\end{exercise}
